% ==============================================================================
% Capítulo 6: Implementação
% ==============================================================================
\chapter{Implementação}
\label{cap:implementacao}

\begin{instrucao}
\orientacao{O que vai na Implementação}
Este capítulo mostra como o projeto foi transformado em software. Inclua:
\begin{itemize}
  \item o ambiente e as ferramentas de desenvolvimento;
  \item a organização do código e os padrões de documentação e codificação
        adotados;
  \item trechos de código que mostrem soluções relevantes ou não óbvias, e
        não o código inteiro;
  \item as telas principais do sistema, relacionadas aos requisitos que
        atendem;
  \item como o sistema foi implantado.
\end{itemize}
O texto não deve ser um manual do usuário. Mostre as telas para evidenciar que
os requisitos foram atendidos e explique o que cada uma permite fazer. Um
manual completo, se existir, vai para um apêndice. Informe o endereço do
repositório com o código-fonte.
\end{instrucao}

Este capítulo descreve a implementação do Monitora. A
\cref{sec:implementacao:organizacao} apresenta a organização do código, a
\cref{sec:implementacao:reserva} detalha a implementação da reserva de
horários, a \cref{sec:implementacao:implantacao} descreve a implantação e a
\cref{sec:implementacao:ia} relata o uso de ferramentas de inteligência
artificial.

% ------------------------------------------------------------------------------
\section{Ambiente e organização do código}
\label{sec:implementacao:organizacao}

O sistema foi desenvolvido no Visual Studio Code, com controle de versões no
Git e repositório hospedado no GitHub\footnote{O endereço do repositório deve
ser informado aqui, por exemplo, \url{https://github.com/usuario/monitora}.}.
O código segue a estrutura padrão de projetos Django, com um diretório por
módulo (\texttt{contas}, \texttt{agenda} e \texttt{relatorios}), cada um com
seus modelos, visões, \emph{templates} e testes. A formatação segue o guia
PEP~8 \cite{vanrossum2001pep8}, verificada automaticamente pela ferramenta
Ruff a cada \emph{commit}. Funções e classes públicas são documentadas com
\emph{docstrings}.

% ------------------------------------------------------------------------------
\section{Reserva de horários}
\label{sec:implementacao:reserva}

O \cref{cod:conflito} mostra o modelo \texttt{Reserva} e a função que cria
uma reserva. A restrição de unicidade declarada nas linhas 13 a 17 implementa
a decisão de projeto discutida na \cref{sec:projeto:decisoes}. Se duas
requisições tentarem reservar o mesmo horário na mesma data, o banco de dados
aceita apenas a primeira e a segunda gera um \texttt{IntegrityError}. A
função \texttt{reservar} (linhas 24 a 33) converte esse erro em uma exceção
do domínio, tratada pela visão, que exibe ao estudante a mensagem de que o
horário acabou de ser ocupado.

\begin{codigo}[htbp]
  \caption{Modelo \texttt{Reserva} e função de reserva com tratamento de conflito}
  \label{cod:conflito}
  \lstinputlisting[language=Python]{codigos/reservas.py}
  \fonteautor
\end{codigo}

\begin{instrucao}
\orientacao{Códigos-fonte}
Inclua apenas trechos curtos e relevantes, em geral com menos de 30 linhas.
O código completo fica no repositório. Coloque o trecho no ambiente
\verb|codigo|, que tem legenda acima, fonte abaixo e lista própria (Lista de
códigos). Para não copiar o código para dentro do \verb|.tex|, guarde o
arquivo na pasta \verb|codigos/| e use \verb|\lstinputlisting|. As opções
\verb|firstline| e \verb|lastline| permitem mostrar só parte do arquivo.
\begin{exemplocodigo}
\begin{codigo}[htbp]
  \caption{Função de reserva}
  \label{cod:reserva}
  \lstinputlisting[language=Python, firstline=24, lastline=33]{codigos/reservas.py}
  \fonteautor
\end{codigo}
\end{exemplocodigo}
Para trechos muito curtos, escritos direto no texto, use o ambiente
\verb|lstlisting| dentro do \verb|codigo|. Explique o código no texto citando
os números das linhas, como no parágrafo acima. Nunca use capturas de tela de
código, que não podem ser lidas por leitores de tela nem copiadas.

Algoritmos devem ser apresentados em pseudocódigo, e não em uma linguagem
específica. O modelo científico mostra como fazer isso com o ambiente
\verb|algoritmo| e o pacote \verb|algpseudocode|.
\end{instrucao}

% ------------------------------------------------------------------------------
\section{Telas do sistema}
\label{sec:implementacao:telas}

\begin{instrucao}
\orientacao{Telas}
Mostre as telas que correspondem aos requisitos principais, uma por figura,
com legenda que diga o que a tela permite fazer, e cite no texto os
requisitos atendidos. Capture as telas em PNG, na resolução original e com
zoom de 100\%, preenchidas com dados fictícios (nunca dados pessoais reais de
estudantes). Para incluir uma captura de tela:
\begin{exemplocodigo}
\begin{figure}[htbp]
  \centering
  \caption{Tela de reserva de horários (RF02 e RF03)}
  \label{fig:tela-reserva}
  \includegraphics[width=0.9\textwidth]{figuras/tecnologico/tela-reserva.png}
  \fonteautor
\end{figure}
\end{exemplocodigo}
Duas telas relacionadas podem ser apresentadas lado a lado com o pacote
\verb|subcaption|, usando \verb|\begin{subfigure}{0.48\textwidth}| dentro do
ambiente \verb|figure|.
\end{instrucao}

A tela de reserva segue o protótipo apresentado na \cref{fig:prototipo}. O
estudante escolhe a disciplina no menu lateral e vê a grade semanal com os
horários livres e o nome de cada monitor (RF02). Ao selecionar um horário, uma
janela de confirmação mostra a data, o monitor e o local do atendimento
(RF03). As reservas do estudante ficam disponíveis no item ``Minhas
reservas'', no qual também é possível cancelá-las (RF04).

% ------------------------------------------------------------------------------
\section{Implantação}
\label{sec:implementacao:implantacao}

O sistema é distribuído em três contêineres Docker, descritos em um arquivo
do Docker Compose, que executam o Nginx, a aplicação Django com o Gunicorn e
o PostgreSQL. A tarefa de envio de lembretes é executada
diariamente pelo agendador do próprio servidor. A versão avaliada no
\cref{cap:testes} foi implantada em uma máquina virtual de homologação
cedida pelo departamento.

% ------------------------------------------------------------------------------
\section{Uso de ferramentas de inteligência artificial}
\label{sec:implementacao:ia}

Durante a implementação, o GitHub Copilot \cite{github2026copilot}, usado
como extensão do Visual Studio Code com o modelo padrão do plano gratuito para
estudantes, sugeriu complementos de código em funções de formatação dos
relatórios e em testes automatizados. Todas as sugestões aceitas foram
revisadas, ajustadas pelo autor e cobertas por testes. Nenhum dado pessoal,
nenhuma credencial e nenhum dado dos usuários da avaliação foram expostos à
ferramenta. A declaração completa e o detalhamento das interações estão no
\cref{ap:ia}.

\begin{instrucao}
\orientacao{Declaração de uso de inteligência artificial}
O Art. 19 da Portaria Normativa nº 310/2026 da Ufes \cite{ufes2026portaria}
torna obrigatória a declaração do uso de \ac{ia} em qualquer produção
acadêmica, exceto os usos meramente acessórios definidos pela comissão de
governança de IA da Ufes (CPGIA). Assistentes de programação, como o GitHub
Copilot, também precisam ser declarados. Para trabalhos acadêmicos, a portaria
pede três coisas. O modelo atende a cada uma delas em um lugar:
\begin{enumerate}
  \item uma declaração com ferramenta, versão, finalidade e forma de
        contribuição, em seção ou apêndice específico (§ 1º): o
        \cref{ap:ia}, que fica como primeiro apêndice para ser encontrado com
        facilidade;
  \item o detalhamento das interações, das motivações e dos métodos de uso
        (§ 2º): também no \cref{ap:ia};
  \item a citação da ferramenta conforme as normas de citação e a indicação,
        na descrição dos métodos, da versão, dos parâmetros e das bases de
        dados usadas (§ 3º): esta seção, com a ferramenta nas Referências.
\end{enumerate}
A portaria também proíbe atribuir autoria ou coautoria a uma ferramenta de IA
(Art. 17) e inserir nela dados pessoais, credenciais institucionais ou dados
confidenciais. Você continua responsável pela correção e pela originalidade
de tudo o que entregar, inclusive do código (Art. 18). Se não usou nenhuma
ferramenta de IA, apague esta seção e o \cref{ap:ia}. Quando o uso decorrer de
tecnologia assistiva, a declaração pode ser apresentada ao setor de apoio à
inclusão, e não no corpo do trabalho (§ 4º). O modelo oficial de declaração
ainda será publicado pela CPGIA. Quando isso acontecer, ajuste o apêndice a
ele.
\end{instrucao}
